A problemática central deste trabalho consiste em verificar se modelos de linguagem generalistas, amplamente utilizados em áreas como a medicina e na educação, são capazes de atender às necessidades específicas da pecuária leiteira, fornecendo respostas úteis e coerentes para auxiliar produtores em processos de tomada de decisão. Para investigar essa questão, diferentes LLMs foram comparados, buscando identificar o desempenho relativo de cada um frente ao domínio estudado.
A avaliação foi conduzida por meio de métricas consolidadas, como BLEU e ROUGE, que medem similaridade textual, e BERTScore e BARTScore, que avaliam proximidade semântica. Adicionalmente, foi utilizado o GPTScore para complementar a análise com uma perspectiva de qualidade contextual. Essa abordagem permite não apenas identificar qual modelo apresentou melhor desempenho, mas também revelar limitações comuns entre eles, apontando a necessidade de adaptações, como RAG, em aplicações futuras.   

\textit{Palavras-Chave}: Large Language Model, Análise, Pecuária Leiteira